\section{Discussion}
\label{sec:Discussion}
The first measurement was the current-voltage measurement. The result is shown in \autoref{fig:IV}. In this plot, the current still rises with larger voltages. When the sensor is fully depleted the gradient of this plot is $m_{\text{IV}} \approx \SI{0.0051}{{\micro\ampere}\per\volt}$. There is also a limit of the linearity, which could causes because of a defect of the control unit. The bias is set to  $\SI{90}{\volt}$. The increasing current is physically plausible, but the possible defect of the control unit limits the validity of the exact numerical value. Therefore, the IV measurement can mainly be used as a qualitative indication of the sensor behaviour.\\\\
The next measurement shows the pedestal signal and the common mode after subtracting the pedestal from the signal. This is shown in \autoref{fig:pedestal}. The pedestal run determines the channel noise without a source, while the common mode measurement checks for a common shift of all channels. This results in the standard deviation $\sigma \approx \num{2.5} \text{ADC Counts}$. Since this noise is small compared to the measured signal amplitudes, the pedestal and common mode correction can be considered reliable for the following measurements.\\\\
The delay of $\SI{64}{\nano\second}$ from the calibration part. This is as expected and similar to the second measured delay of $\SI{61}{\nano\second}$. This measurement is shown in \autoref{fig:delay}. The agreement of both delay values supports the validity of the timing calibration.\\
There is also the result of the calibration for five different strips. The last one seems unphysical. This could results from a system restart with wrong settings of the calibration run. The other four measurements are as expected and leads to the mean fit curve with the parameters given in \autoref{tab:calibration_parameters}. The calibration measurement is given in \autoref{fig:calibration}. There is also one measurement with $\SI{0}{\volt}$ which is shown in \autoref{fig:calibration_0}. This shows that with a bias of $\SI{0}{\volt}$ the measured ADC Counts are lower compared to the previously measured values. Since the calibration charge is injected directly into the strips, this is probably caused by additional detector noise and background at working voltage, which is not present at $\SI{0}{\volt}$. Because one calibration run is unphysical, it should not be weighted equally with the other measurements. The remaining calibration points still show a consistent trend, so the calibration is usable. \\\\
The next step was to analyze the stip structure within the laser. The results are shown in \autoref{fig:laser}. The strip distance is estimated from the local minima of the summed signal, because these minima correspond to positions where the laser is above a strip. With the assumed step size this gives $d_{\text{strip}}=\left(\num{27.3} \pm \num{4.7}\right)\si{\micro\meter}$. This value is too small for the strip distance, so the micrometer was probably changed in a different way than assumed during the measurement. If the micrometer was moved by $\SI{100}{\micro\meter}$ instead of $\SI{10}{\micro\meter}$ per measurement, the distance has to be scaled to $d_{\text{strip}}=\left(\num{273} \pm \num{47}\right)\si{\micro\meter}$. This value is only a rough estimate and is larger than expected, so the minima are not resolved well enough for a reliable strip distance. Further, this measurement confirms the strip structure qualitatively, but the exact strip distance, strip width and laser width cannot be determined reliably from this scan. \\\\
The last measurement was about the charge collection efficiency. The result is given in \autoref{fig:CCE}. With a fit the penetration depth $a = \left(\num{323}\pm\num{36}\right)\si{\micro\meter}$, the scaling factor $A= \num{0.945}\pm\num{0.008}$ and the gradient of the plateau $m=\left(\num{0.0006}\pm\num{0.0001}\right)\si{\per\volt}$ is calculated. Since this penetration depth is slightly larger than the detector thickness of $\SI{300}{\micro\meter}$, the fit is at the limit of the physically reasonable region. A fit only to the increasing part up to $\SI{80}{\volt}$ gives $a=\left(\num{309}\pm\num{58}\right)\si{\micro\meter}$, which agrees better with the detector thickness. The fitted plateau and the small gradient are consistent with the expectation that the charge collection becomes almost constant after the depletion. This supports the validity of the charge collection efficiency measurement, although systematic effects from the calibration and electronics can still influence the absolute values.\\\\
In summary this protocol introduces the work with Silicon Strip detectors well.
The main detector properties could be observed in the different measurements, especially the influence of the bias voltage, the electronic noise and the charge collection behaviour.
Some measurements are limited by systematic effects, such as the unphysical calibration point, the possible wrong micrometer step size in the laser scan and the uncertainty of the CCE fit.
For this reason, the experiment gives a valid qualitative characterization of the sensor.
